\section*{A Proof of Theorem 13 (author source pp. 13:11--13:12)}
The proof uses known results on quantifier elimination which we first outline. We follow the monograph of Basu, Pollack and Roy [1]. Theorem 13 requires an elimination of only a single block of existential quantifiers, so we focus on this case only.
For $b\in\RR$, let
\[\sgn(b):=\begin{cases}1,&b>0,\\0,&b=0,\\-1,&b<0.\end{cases}\]
Given $b=\langle b_1,\ldots,b_m\rangle\in\RR^m$, let $\sgn(b):=\langle\sgn(b_1),\ldots,\sgn(b_m)\rangle\in\{-1,0,1\}^m$. Let $F=F(z,y)$ be a sequence of $m$ polynomials $f_1,\ldots,f_m\in\RR[z,y]$ in variables $z=\{z_1,\ldots,z_{k_1}\}$ and $y=\{y_1,\ldots,y_{k_2}\}$. Given $a\in\RR^{k_1}$, define $\SGN_1(F,a)\subseteq\{-1,0,1\}^m$
\[\SGN_1(F,a):=\{\sgn(F(a,b)):b\in\RR^{k_2}\}.\]
Let
\[\SGN(F):=\{\SGN_1(F,a):a\in\RR^{k_1}\}.\]
Theorem 14.16 from [1] provides the following bound on the size of SGN:
\begin{originalresult}{Theorem ([1])}
If every polynomial in $F$ has degree at most $d$ then
\[|\SGN(F)|\le m^{(k_1+1)(k_2+1)}d^{O(k_1)O(k_2)}.\tag{4}\]
\end{originalresult}
We now apply this result to the case of Theorem 13. Let $V,s,\mathcal A$ be as in the assumption. Every $A\in\mathcal A$ can be described by a linear system with $s$ inequalities. Namely, for every $x\in V$,
\[x\in A\text{ iff }\exists y\in\RR^s\ Cx+Dy=b,\ y\ge0,\tag{5}\]
where $C\in\RR^{t\times r}$, $D\in\RR^{t\times s}$ and $b\in\RR^t$. Since $Cx+Dy=b$ is a system of equations in $r+s$ variables $x,y$, we can also assume $t=r+s$.
Let us view the parameters $C,D,b$ in (5) as variables. Let $z$ be the set of these variables, of size $k_1=(r+s)(r+s+1)$. Given $v\in V$, let $F_v(z,y)$ be the sequence of $(r+s)$ polynomials
\[Cv+Dy-b\]
in variables $z$ and $y=\{y_1,\ldots,y_s\}$. Let $F(z,y)$ be the union of $F_v(z,y)$ over all $v\in V$, together with the polynomials $y_1,\ldots,y_s$. Hence $F$ consists of $m=s+|V|(r+s)$ polynomials of degree at most two.
$F(z,y)$ is set up so that
\[|\mathcal A|\le|\SGN(F)|.\]
To see this, observe that whenever the parameters $z$ are fixed, the set $A\subseteq V$ given by (5) is uniquely determined by $\SGN_1(F(z,y))$. Since every $A\in\mathcal A$ is obtained by some fixing of the parameters, we indeed obtain $|\mathcal A|\le|\SGN(F(z,y))|$.
Finally, we can apply (4) to estimate $|\SGN(F)|$ with $m=s+|V|(r+s)$, $k_1=(r+s)(r+s+1)$, $k_2=s$, and $d=2$. To simplify the expression, we can assume $s+r\le|V|$; otherwise the upper bound asserted in Theorem 13 exceeds the trivial bound $|\mathcal A|\le2^{|V|}$. This means that $m\le2|V|^2$. If we loosen the bound (4) as $|\SGN(F)|\le(dm)^{O(k_1)O(k_2)}$, we obtain (recall that $s\ge1$)
\[|\SGN(F)|\le2^{O(s(s+r)^2\log|V|)},\]
as required.
